\documentclass[main.tex]{subfiles}
\begin{document}
\section{Minimizer-space de Bruijn graphs: Whole-genome assembly of long reads in minutes on a personal computer.}
\subsection{Eckdaten}
\begin{itemize}
    \item \textbf{Länge:} 20-45 min
\end{itemize}  
Ekim, Barış, Bonnie Berger, and Rayan Chikhi. “Minimizer-space de Bruijn graphs: Whole-genome assembly of long reads in minutes on a personal computer.” Cell systems 12.10 (\textbf{2021}): 958-968. 
\subsection{Erarbeitung}
\subsubsection{Was macht der überhaupt?} 
\begin{itemize}
    \item Whole-genome assembly 
\end{itemize}
\subsubsection{Wieso wurde der entwickelt?}  
\begin{itemize}
    \item DNA-Sequenzierung entwickelt sich weiter und die daraus entstehenden Reads werden länger und die zusammenzusetzen ist schwierig 
    \begin{itemize}
        \item DNA sequencing data continue to progress toward longer reads with increasingly lower sequencing error rates
        \item High-throughput DNA sequencing
    \end{itemize}
    \item improvements in speed and memory usage for human genome assembly and metagenome assembly 
    \begin{itemize}
        \item achieves orders-of-magnitude improvement in both speed and memory usage over existing methods without compromising accuracy
        \item It improves the efficiency of \gls{v:pgc}
        \item here has traditionally been a trade-off between algorithmic efficiency and loss of information, at least during the initial sequence-processing steps
    \end{itemize}
\end{itemize} 
\subsection{Was ist neu?}
\begin{itemize}
    \item \gls{v:mssa}
    \item expanding the alphabet of DNA sequences to \gls{v:atomictokens} made of fixed-length words
    \item  instead of building an assembly over sequence bases—the standard approach that for clarity we refer to as base space—newly performs assembly in \gls{v:minimizer space} and later converts it back to base-space assemblies 
    \item Specifically, each read is initially converted to an ordered sequence of its \glspl{v:minimizer}
\end{itemize}

\subsubsection{Welche Datenstrukturen und Strategien benutzt er?}
\subsubsection{Was macht er besser als andere?}
\subsubsection{Was macht er schlechter als andere?}

\subsubsection{Aufgekommene Fragen}
\begin{itemize}
    \item Wie ist es möglich aus Minimizern die ursprüngliche Sequenz wieder zu rekonstruieren? 
    \item Werden Window-Minimizer oder globale Minimizer verwendet? 
    
    \item wie werden die Minimizer ausgesucht? 
    
    \item wie wird festgestellt, ob eine Sequenz high oder low errror ist? 
\end{itemize}
 
\end{document}